% ============================================================================
% preamble.tex -- all machinery, one edition (Polish).  House colours are the
% estate's (research-documentation).
% CONTRACT: main.tex defines \bookstatus, \bookpdftitle BEFORE it \inputs this.
% ============================================================================
\usepackage[T1]{fontenc}
\usepackage[utf8]{inputenc}
\usepackage{lmodern}
\usepackage[main=polish]{babel}
\usepackage[a4paper,left=2.4cm,right=2.4cm,top=2.6cm,bottom=2.6cm]{geometry}
\usepackage{microtype}
\usepackage{xcolor}
\usepackage{etoolbox}
\usepackage{enumitem}
\usepackage{booktabs}
\usepackage{longtable}
\usepackage{array}
\usepackage{colortbl}
\usepackage{multicol}
\usepackage{fancyhdr}
\usepackage{titlesec}
\usepackage[most]{tcolorbox}
\usepackage{amssymb}

\definecolor{primaryblue}{RGB}{41,128,185}
\definecolor{darkblue}{RGB}{31,78,121}
\definecolor{lightgray}{RGB}{236,240,241}
\definecolor{darkgray}{RGB}{52,73,94}
\definecolor{accentgreen}{RGB}{39,174,96}
\definecolor{accentorange}{RGB}{230,126,34}
\definecolor{accentred}{RGB}{231,76,60}

\setlength{\parindent}{0pt}
\setlength{\parskip}{0.55em plus 0.2em minus 0.1em}
\setlist{nosep,leftmargin=1.6em}
\renewcommand{\arraystretch}{1.12}

% ---- computed values: \val{key} reads generated/values.tex ----------------
% GENERATED by tools/gen.py -- do not edit
\csgdef{gval@bread.1.hydration}{70}
\csgdef{gval@bread.1.mass}{862}
\csgdef{gval@bread.2.hydration}{70}
\csgdef{gval@bread.2.mass}{860}
\csgdef{gval@bread.3.hydration}{82,6}
\csgdef{gval@bread.3.mass}{1060}
\csgdef{gval@bread.4.hydration}{80}
\csgdef{gval@bread.4.mass}{913}
\csgdef{gval@bread.5.hydration}{60}
\csgdef{gval@bread.5.mass}{813}
\csgdef{gval@bread.6.hydration}{64}
\csgdef{gval@bread.6.mass}{835}
\csgdef{gval@bread.7.hydration}{70}
\csgdef{gval@bread.7.mass}{861}
\csgdef{gval@bread.8.hydration}{85}
\csgdef{gval@bread.8.mass}{937}
\csgdef{gval@bread.hydration.max}{85}
\csgdef{gval@bread.hydration.min}{60}
\csgdef{gval@breads.total}{8}
\csgdef{gval@chapters.total}{10}
\csgdef{gval@exercises.total}{73}
\csgdef{gval@tasks.total}{10}
\csgdef{gval@traps.total}{11}
\csgdef{gval@verify.total}{55}
\csgdef{gval@week.A.free}{0,5}
\csgdef{gval@week.A.sleep}{12}
\csgdef{gval@week.B.free}{0,25}
\csgdef{gval@week.B.sleep}{12}
\csgdef{gval@week.C.free}{5}
\csgdef{gval@week.C.sleep}{12}
\csgdef{gval@week.templates}{3}
\csgdef{gval@written.total}{10}

% GENERATED by tools/gen.py -- do not edit
\csgdef{chtitle@01}{Budżet doby}
\csgdef{chlead@01}{Doba ma 24 godziny: budżet czasu, w którym 12 godzin zajmuje sen}
\csgdef{chtitle@02}{Ile snu potrzebujesz}
\csgdef{chlead@02}{Uczciwie: ile snu potrzebuje dorosły, kiedy 12 godzin jest sygnałem do lekarza, a kiedy wyborem}
\csgdef{chtitle@03}{Co wyciąć}
\csgdef{chlead@03}{Audyt czasu: co realnie zjada dobę i jak odzyskać godziny bez heroizmu}
\csgdef{chtitle@04}{Praca w osiem godzin}
\csgdef{chlead@04}{Granice pracy, dojazd, asynchroniczność i wieczór, który się nie rozlewa}
\csgdef{chtitle@05}{Rytm snu}
\csgdef{chlead@05}{Stała pora, światło, kofeina, sypialnia i wieczorna rutyna}
\csgdef{chtitle@06}{Szczelny dzień}
\csgdef{chlead@06}{Bloki czasu, szablony dnia i tygodnia, które się sumują do 24 godzin}
\csgdef{chtitle@07}{Dom w minimum}
\csgdef{chlead@07}{Zakupy, gotowanie hurtowe, sprzątanie w progu i automatyzacja domowego życia}
\csgdef{chtitle@08}{Chleb robi się sam}
\csgdef{chlead@08}{Dlaczego chleb pasuje do długiego snu: fermentacja na zimno i czas bierny}
\csgdef{chtitle@09}{Piekarnia w rytmie snu}
\csgdef{chlead@09}{Harmonogramy wypieków godzina po godzinie, kilka rodzajów chleba w tygodniu}
\csgdef{chtitle@10}{Tydzień, podróż, kryzys}
\csgdef{chlead@10}{Utrzymanie rytmu: dzieci, podróże, choroba, awaria i wracanie do planu}

\newcommand{\val}[1]{\ifcsname gval@#1\endcsname\csuse{gval@#1}\else\textbf{??}\fi}
\newcommand{\twodig}[1]{\ifnum#1<10 0\fi\number#1}
\newcommand{\code}[1]{\texttt{#1}}
% a claim that needs a primary source or a specialist before you rely on it
\newcommand{\verify}[1]{#1\textsuperscript{\textcolor{accentorange}{[?]}}}

\addto\captionspolish{%
  \renewcommand{\chaptername}{Rozdział}%
  \renewcommand{\appendixname}{Dodatek}%
  \renewcommand{\contentsname}{Spis treści}}

\titleformat{\part}[display]
  {\centering\color{primaryblue}\sffamily\bfseries\Huge}{}{0pt}{}
\titleformat{\chapter}[display]
  {\color{darkblue}\sffamily\bfseries}
  {\Large\color{primaryblue}\MakeUppercase{\chaptertitlename}\ \thechapter}{0.4em}{\Huge}[\vspace{0.3em}\titlerule]
\titlespacing*{\chapter}{0pt}{-10pt}{22pt}
\titleformat{\section}{\color{primaryblue}\Large\bfseries\sffamily}{\thesection}{0.6em}{}
\titleformat{\subsection}{\color{darkblue}\large\bfseries\sffamily}{}{0pt}{}
\titlespacing*{\subsection}{0pt}{1.1em}{0.3em}

\pagestyle{fancy}
\fancyhf{}
\renewcommand{\headrulewidth}{0.4pt}
\fancyhead[L]{\small\sffamily\color{darkgray}\nouppercase{\leftmark}}
\fancyhead[R]{\small\sffamily\color{accentred}\bookstatus}
\fancyfoot[C]{\small\thepage}
\fancypagestyle{plain}{\fancyhf{}\renewcommand{\headrulewidth}{0pt}\fancyfoot[C]{\small\thepage}}

% ---- boxes ------------------------------------------------------------------
\tcbset{mfa/.style={enhanced,breakable,boxrule=0pt,frame hidden,left=9pt,right=8pt,top=5pt,bottom=5pt,
  before skip=9pt,after skip=9pt,borderline west={3pt}{0pt}{#1},colback=#1!7}}
\newcommand{\boxlead}[2]{{\sffamily\bfseries\color{#1}#2.}\ }
\newtcolorbox{falle}[1][Pułapka]{mfa=accentred,before upper={\boxlead{accentred}{#1}}}
\newtcolorbox{zasada}[1][Zasada]{mfa=accentgreen,before upper={\boxlead{accentgreen}{#1}}}
\newtcolorbox{uwaga}[1][Uwaga]{mfa=darkgray,before upper={\boxlead{darkgray}{#1}}}
\newtcolorbox{piekarnik}[1][Z piekarnika]{mfa=accentorange,before upper={\boxlead{accentorange}{#1}}}
\newtcolorbox{przyklad}[1][Przykład]{mfa=primaryblue,before upper={\boxlead{primaryblue}{#1}}}

\newenvironment{cele}
  {\begin{tcolorbox}[mfa=primaryblue]{\sffamily\bfseries\color{primaryblue}Po tym rozdziale potrafisz:}\begin{itemize}[label=$\square$]}
  {\end{itemize}\end{tcolorbox}}

\newcommand{\zadanie}[1]{%
  \begin{tcolorbox}[mfa=accentgreen]
  {\sffamily\bfseries\color{accentgreen}Zadanie na ten tydzień (zrób, nie tylko przeczytaj).}\ #1
  \end{tcolorbox}}

\newcommand{\chapteropen}[1]{%
  \chapter[\csuse{chtitle@#1}]{\csuse{chtitle@#1}}%
  \begin{center}\small\sffamily\color{darkgray}\csuse{chlead@#1}\end{center}}

% ---- generated tables ---------------------------------------------------------
\newcommand{\generated}[1]{\input{generated/#1}}

% ---- exercises: answers are written where the exercise is, printed at the back
\newcommand{\answerstore}{}
\newcommand{\lastansch}{0}
\newenvironment{cwiczenia}
  {\begin{enumerate}[label=\textbf{\arabic*.},leftmargin=2em,itemsep=0.35em]}
  {\end{enumerate}}
\newcommand{\exitem}[2]{%
  \item #1%
  \begingroup\edef\tmp{\endgroup\noexpand\gappto\noexpand\answerstore{%
    \noexpand\answeritem{\number\value{chapter}}{\number\value{enumi}}{\unexpanded{#2}}}}\tmp}
\newcommand{\answeritem}[3]{%
  \ifnum#1=\lastansch\else
    \gdef\lastansch{#1}\par\medskip\noindent{\sffamily\bfseries\color{primaryblue}Rozdział #1 --- \csname chtitle@\twodig{#1}\endcsname}\par\nobreak
  \fi
  \noindent\textbf{#1.#2}\ #3\par}

\usepackage[colorlinks=true,linkcolor=primaryblue,urlcolor=primaryblue,citecolor=primaryblue,
  pdftitle={\bookpdftitle},pdfauthor={Konrad Cinkusz},pdflang={pl-PL},bookmarksnumbered=true]{hyperref}
